% =====================================================================
%  packages.tex  --  preamble: packages, fonts, sizes, page layout
%  Engine: LuaLaTeX (required for fontspec + Times New Roman + microtype
%  letterspacing).  See CLAUDE.md for the full style specification.
% =====================================================================

% ---------------------------------------------------------------------
%  Math (load amsmath BEFORE unicode-math) and fonts
% ---------------------------------------------------------------------
\usepackage{amsmath}

\usepackage{unicode-math}          % loads fontspec; OpenType math support
\setmainfont{Times New Roman}      % body font = real Times New Roman
\setmathfont{TeX Gyre Termes Math} % Times-clone math font (closest match)

\usepackage{microtype}             % micro-typography + \textls letterspacing

% ---------------------------------------------------------------------
%  Exact font sizes  (base class is 11pt; redefine the size macros so
%  they hit the values requested verbatim)
%     \normalsize = 11pt   \large = 12pt   \Large = 14pt   \LARGE = 16pt
% ---------------------------------------------------------------------
\makeatletter
\renewcommand\normalsize{%
  \@setfontsize\normalsize{11}{13.6}%
  \abovedisplayskip 11\p@ \@plus3\p@ \@minus6\p@
  \abovedisplayshortskip \z@ \@plus3\p@
  \belowdisplayshortskip 6.5\p@ \@plus3.5\p@ \@minus3\p@
  \belowdisplayskip \abovedisplayskip}
\renewcommand\large{\@setfontsize\large{12}{14.5}}
\renewcommand\Large{\@setfontsize\Large{14}{17}}
\renewcommand\LARGE{\@setfontsize\LARGE{16}{19}}
\makeatother
\normalsize                        % apply the new \normalsize now

% ---------------------------------------------------------------------
%  Page geometry  --  2.5 cm margins all round.
%  No includehead/includefoot: the ruled header & footer sit in the top and
%  bottom margins (Word-style), so the 2.5 cm text block keeps its full size
%  (the head/sep/foot are not subtracted from the body).
% ---------------------------------------------------------------------
\usepackage[a4paper,margin=2.5cm,
            headheight=15pt,headsep=9pt,footskip=35pt]{geometry}
%% headsep/footskip tuned so the header rule sits ~21.8 mm from the top and the
%% footer rule ~17.8 mm from the bottom, matching HW_documentation_word.pdf
%% (the 2.5 cm body margins are unchanged, so the text block keeps its size).

% ---------------------------------------------------------------------
%  Paragraph layout  --  block style, 8pt between paragraphs (and titles)
% ---------------------------------------------------------------------
\setlength{\parindent}{0pt}
\setlength{\parskip}{8pt plus 2pt minus 1pt}

% ---------------------------------------------------------------------
%  Heading styles (titlesec).  Word-style multilevel list: the NUMBER sits at
%  the left margin and only the TITLE is indented past it (the number is boxed
%  to the indent width).  Unnumbered headings (Appendix, References, Contents)
%  have no number, so their title sits at the left margin -> no indentation.
%     "chapter"    -> \section      "1.  Title"     14pt bold, title @ 0.75cm
%     "section"    -> \subsection   "2.1 Title"     12pt bold, title @ 0.75cm
%     "subsection" -> \subsubsection "2.1.1 Title"  12pt      , title @ 1.25cm
%  Top level shows a trailing period (1.) like the Word doc.
% ---------------------------------------------------------------------
\usepackage{titlesec}

\titleformat{\section}[hang]
  {\normalfont\Large\bfseries}{\makebox[0.75cm][l]{\thesection.}}{0pt}{}
\titlespacing*{\section}{0pt}{8pt}{8pt}

\titleformat{\subsection}[hang]
  {\normalfont\large\bfseries}{\makebox[0.75cm][l]{\thesubsection}}{0pt}{}
\titlespacing*{\subsection}{0pt}{8pt}{8pt}

\titleformat{\subsubsection}[hang]
  {\normalfont\large}{\makebox[1.25cm][l]{\thesubsubsection}}{0pt}{}
\titlespacing*{\subsubsection}{0pt}{8pt}{8pt}

% Number down to subsubsection and list them in the ToC
\setcounter{secnumdepth}{3}
\setcounter{tocdepth}{3}

% In the ToC, give the top-level (section) number a trailing period -> "1."
% (matching the body headings, where titlesec prints "\thesection.").  Deeper
% levels (2.1, 2.1.1) keep the bare number, and the unnumbered Appendix entry
% has no \numberline so it is unaffected.  This is a copy of article's
% \l@section with a period injected into the number box.
\makeatletter
\renewcommand*\l@section[2]{%
  \ifnum \c@tocdepth >\z@
    \addpenalty\@secpenalty
    \addvspace{1.0em \@plus\p@}%
    \setlength\@tempdima{1.5em}%
    \begingroup
      \parindent \z@ \rightskip \@pnumwidth
      \parfillskip -\@pnumwidth
      \leavevmode \bfseries
      \advance\leftskip\@tempdima
      \hskip -\leftskip
      \def\numberline##1{\hb@xt@\@tempdima{##1.\hfil}}%
      #1\nobreak\hfil \nobreak\hb@xt@\@pnumwidth{\hss #2}\par
    \endgroup
  \fi}
\makeatother

% ---------------------------------------------------------------------
%  Submission date  ->  "YYYY. MM. DD." (auto from compile date).
%  Not used by default; available as \docdate (e.g. on the title page).
% ---------------------------------------------------------------------
\newcommand{\twodig}[1]{\ifnum#1<10 0\fi\number#1}
\newcommand{\docdate}{\number\year.\ \twodig{\month}.\ \twodig{\day}.}

% ---------------------------------------------------------------------
%  Ruled header & footer (fancyhdr).  Footer text is \normalsize (11pt).
%     header : top rule only (no text), ~2 blank lines down to the body
%     footer : page number (right) above a rule
%     page numbering starts on the title page (counted as 1) but is only
%     shown from page 2 onward -> the title page uses an empty-text style
% ---------------------------------------------------------------------
\usepackage{fancyhdr}

\fancypagestyle{main}{%
  \fancyhf{}%
  \renewcommand{\headrulewidth}{0.4pt}%
  \renewcommand{\footrulewidth}{0.4pt}%
  \fancyfoot[R]{\normalsize\thepage}%
}

\fancypagestyle{titlepg}{%
  \fancyhf{}%
  \renewcommand{\headrulewidth}{0.4pt}%
  \renewcommand{\footrulewidth}{0.4pt}%
  % invisible page-number-height strut so the footer rule sits at the same
  % height as on numbered pages (where the page number gives the foot its height)
  \fancyfoot[R]{\normalsize\vphantom{\thepage}}%
}

\pagestyle{main}

% ---------------------------------------------------------------------
%  Figures, tables, captions, code listings
% ---------------------------------------------------------------------
\usepackage{graphicx}
\usepackage{xcolor}                 % colours for the listing frame/background
\usepackage{float}
\usepackage{booktabs}
\usepackage{multirow}              % material name/density spanning component rows
\usepackage{tabularx}              % equal-width columns (ICRP table at fixed width)
% equal-width X-column variants: Y centres the content, L flushes it left
\newcolumntype{Y}{>{\centering\arraybackslash}X}
\newcolumntype{L}{>{\raggedright\arraybackslash}X}
\usepackage{caption}
\captionsetup{font=it}             % italic captions, e.g. "Figure 1: ..."
% table captions sit ABOVE the table (already placed before the tabular) and are
% flush LEFT; singlelinecheck=false keeps short one-line captions left-aligned too.
% Scoped to tables so figure captions keep their centred, below-the-figure style.
\captionsetup[table]{position=above,justification=raggedright,singlelinecheck=false}

\usepackage{pgfplots}              % TikZ-based plotting (ICRP coefficient curve)
\pgfplotsset{compat=1.18}
\usetikzlibrary{arrows.meta}       % arrow tips for the plot axes
\usetikzlibrary{patterns,patterns.meta}  % material hatching (fuel-rod blueprint)

\usepackage{listings}              % for the MCNP input file in the appendix
% Consolas as the listing font; ligatures off so e.g. "->" stays literal
\newfontfamily\lstfont{Consolas}[Ligatures={NoCommon,NoContextual}]
% \ct{RRGGBB}{text}: a Consolas-8pt coloured token, used in the listing's
% escape-to-LaTeX markup.  The appendix MCNP dump (inlined in
% chapters/appendix.tex) was pre-tokenised with the VS Code MCNP grammar so the
% colours match the editor exactly.
\newcommand{\ct}[2]{{\lstfont\fontsize{8}{10}\selectfont\textcolor[HTML]{#1}{#2}}}
% \cb{text}: a bold Consolas-8pt listing token (no colour) -- used for the
% MCNP title card on the first line of the appendix listing.
\newcommand{\cb}[1]{{\lstfont\fontsize{8}{10}\selectfont\textbf{#1}}}
\lstset{
  basicstyle=\lstfont\fontsize{8}{10}\selectfont,
  numbers=left,                    % line numbering on the left
  numberstyle=\lstfont\fontsize{8}{10}\selectfont\color[HTML]{454545},
  numbersep=8pt,
  stepnumber=1,
  frame=single,                    % border around the listing
  backgroundcolor=\color[gray]{0.95}, % light-gray background
  escapeinside={(*@}{@*)},         % inline colour markup from the colouriser
  breaklines=true,
  columns=fullflexible,
  keepspaces=true,
  showstringspaces=false,
  % frame spans 98% of \textwidth, centred; line numbers sit flush to the left border
  framexleftmargin=23pt,
  xleftmargin=\dimexpr0.01\textwidth+23pt\relax,
  xrightmargin=0.01\textwidth,
}

% ---------------------------------------------------------------------
%  Bibliography (biblatex + biber).  refs.bib lives in bib/.
% ---------------------------------------------------------------------
\usepackage[backend=biber,style=numeric,sorting=none]{biblatex}
\addbibresource{bib/refs.bib}

% ---------------------------------------------------------------------
%  Links (load hyperref late)
% ---------------------------------------------------------------------
\usepackage[hidelinks]{hyperref}
